%================================================
% INTRODUCTION
%================================================
\section{Giới thiệu}

\subsection{Bối cảnh bài lab}

Lab 03 của môn Nhập môn Dữ liệu lớn yêu cầu sinh viên giải quyết bốn bài toán phân tích dữ liệu phân tán trên bộ dữ liệu \textbf{Amazon Sale Report}, sử dụng hai nhóm công nghệ:

\begin{itemize}
    \item \textbf{Task 1-1 và Task 1-2:} MapReduce nâng cao --- triển khai thuật toán bằng các phép biến đổi Spark RDD theo tư duy xử lý cặp khóa - giá trị.
    \item \textbf{Task 2-1 và Task 2-2:} Spark Structured APIs --- sử dụng DataFrame/Dataset API với bộ tối ưu hóa Catalyst.
\end{itemize}

\subsection{Bộ dữ liệu Amazon Sale Report}

Bộ dữ liệu chứa thông tin đơn hàng trên nền tảng Amazon tại thị trường Ấn Độ, với khoảng \textbf{128,000 dòng} và \textbf{24 cột}. Các cột chính được sử dụng trong bốn task:

\begin{table}[H]
\centering
\begin{tabular}{lll}
\toprule
\textbf{Cột} & \textbf{Kiểu} & \textbf{Mô tả} \\
\midrule
\texttt{Date}              & String  & Ngày đặt hàng, định dạng \texttt{MM-dd-yy} \\
\texttt{Status}            & String  & Trạng thái đơn (Shipped, Cancelled, ...) \\
\texttt{Fulfilment}        & String  & Đơn vị thực hiện (Amazon, Merchant) \\
\texttt{ship-service-level}& String  & Mức dịch vụ (Standard, Expedited) \\
\texttt{ship-city}         & String  & Thành phố giao hàng \\
\texttt{ship-state}        & String  & Bang giao hàng \\
\texttt{SKU}               & String  & Mã sản phẩm \\
\texttt{Style}             & String  & Mã kiểu dáng sản phẩm \\
\texttt{Size}              & String  & Size sản phẩm (S, M, L, XL, XXL, ...) \\
\texttt{Qty}               & Integer & Số lượng mua \\
\texttt{Amount}            & Double  & Giá trị đơn hàng (INR) \\
\texttt{promotion-ids}     & String  & Danh sách mã khuyến mãi \\
\texttt{Courier Status}    & String  & Trạng thái vận chuyển \\
\bottomrule
\end{tabular}
\caption{Các trường dữ liệu chính từ Amazon Sale Report dùng trong báo cáo}
\label{tab:amazon_dataset_columns}
\end{table}

\subsection{Môi trường thực thi}

\begin{table}[H]
\centering
\begin{tabular}{ll}
\toprule
\textbf{Thành phần} & \textbf{Chi tiết} \\
\midrule
Apache Spark  & 3.4.0 \\
Scala         & 2.13.8 \\
Chế độ chạy   & Chế độ chạy cục bộ (\texttt{local[*]}) \\
Hệ điều hành  & Linux chạy trực tiếp trên máy cá nhân \\
Java          & OpenJDK 8 (1.8.0\_482) \\
\bottomrule
\end{tabular}
\caption{Thông số cấu hình môi trường thực thi chương trình}
\label{tab:execution_environment}
\end{table}

\subsection{Lý do chọn Scala}

Scala được chọn làm ngôn ngữ triển khai cho toàn bộ bốn task vì ba lý do chính:
\begin{enumerate}
    \item \textbf{Tối ưu điểm số ngôn ngữ:} Theo thang điểm đánh giá, sử dụng Scala cho giải pháp hoàn chỉnh được cộng tối đa 0.125 điểm cho mỗi task.
    \item \textbf{An toàn kiểu dữ liệu:} Scala là ngôn ngữ định kiểu tĩnh, giúp phát hiện lỗi trong quá trình biên dịch thay vì khi chạy chương trình --- điều này đặc biệt quan trọng khi làm việc với Dataset API.
    \item \textbf{API gốc của Spark:} Spark được viết bằng Scala; các API như \texttt{aggregateByKey}, \texttt{reduceByKey}, hay các bộ mã hóa của Dataset hoạt động tự nhiên nhất với Scala, tránh được chi phí hiệu năng dư thừa do quá trình tuần tự hóa của PySpark.
\end{enumerate}

\subsection{Tổng quan kết quả đầu ra}

\begin{table}[H]
\centering
\begin{tabular}{llll}
\toprule
\textbf{Task} & \textbf{Công nghệ} & \textbf{Tệp đầu ra} & \textbf{Trạng thái} \\
\midrule
Task 1.1 & MapReduce (RDD) & \texttt{Task\_1-1.csv} & Hoàn thành, đã kiểm chứng \\
Task 1.2 & MapReduce (RDD) & \texttt{Task\_1-2.csv} & Hoàn thành, đã kiểm chứng \\
Task 2.1 & Spark Structured API & \texttt{Task\_2-1.parquet} & Hoàn thành, đã kiểm chứng \\
Task 2.2 & Spark Structured API & \texttt{Task\_2-2.parquet} & Hoàn thành, đã kiểm chứng \\
\bottomrule
\end{tabular}
\caption{Tổng quan về các task thực hiện và định dạng tệp đầu ra tương ứng}
\label{tab:tasks_summary}
\end{table}

\subsection{Cấu trúc báo cáo}

Báo cáo được tổ chức thành các phần sau: thông tin nhóm (Chương 1), sau đó mỗi task là một chương riêng biệt (Chương 2--5) với cấu trúc thống nhất gồm: phân tích truy vấn, phân rã bước, lý do phân rã, chiến lược thực thi kèm các đoạn mã nguồn minh họa, phân tích chuyên sâu, kết quả và kiểm chứng. Phần kết luận (Chương 6) tổng kết sự cân nhắc đánh giá giữa MapReduce và Spark Structured APIs.
